\documentclass[11pt]{article}
\usepackage[a4paper,margin=1in]{geometry}
\usepackage{amsmath,amssymb}
\usepackage{hyperref}
\title{Hexagonal and Trigonal Conventions in PyTex}
\author{PyTex Contributors}
\date{\today}

\begin{document}
\maketitle

\section{Purpose}

Hexagonal and trigonal systems are a common source of avoidable confusion because internal coordinate storage, direct-space notation, reciprocal-space notation, and teaching conventions are often mixed without warning.

\section{Canonical Policy}

PyTex fixes the following rules:

\begin{itemize}
  \item Internal direct-space vectors use a three-component direct basis.
  \item Internal reciprocal-space planes use a three-component reciprocal basis.
  \item Four-index notation is an interface and documentation convention, not an internal storage requirement.
\end{itemize}

\section{Direction Mapping}

For three-index direct-basis directions $[u\,v\,w]$, the four-index Weber form follows
\begin{align*}
  U &= \frac{2u - v}{3}, \\
  V &= \frac{2v - u}{3}, \\
  T &= -\frac{u + v}{3}, \\
  W &= w,
\end{align*}
with $U + V + T = 0$.

\section{Plane Mapping}

For reciprocal-space planes, the Miller--Bravais redundancy is
\[
  i = -(h + k),
\]
so the four-index plane $(h\,k\,i\,l)$ is normalized internally to the three-component reciprocal-basis form $(h\,k\,l)$.

\section{Normative References}

\begin{thebibliography}{9}
\bibitem{degraef2003}
M. De Graef,
\textit{Introduction to Conventional Transmission Electron Microscopy},
Cambridge University Press, 2003.
DOI: \url{https://doi.org/10.1017/CBO9780511615092}.

\bibitem{ita}
Th. Hahn (ed.),
\textit{International Tables for Crystallography, Volume A: Space-Group Symmetry},
IUCr / Springer.
DOI: \url{https://doi.org/10.1107/97809553602060000100}.
\end{thebibliography}

\end{document}
